\section{Discussion}

The pre-registered primary endpoint was not met. Added to six published
competitor features under leakage-controlled grouping, SFI left the combined
real-cohort gradient-boosting increment at $\dAUC=+0.012$, with a bootstrap
interval $[-0.017,+0.044]$ that excludes the pre-registered $0.05$ gate as well
as containing zero (Table~1, Figure~1). On $182$ patient-derived anatomies the
effect was measured and came out small; power is not the limiting factor. Small
is not zero. The lower bound of $-0.017$ leaves an increment of a few
thousandths compatible with these data, and only a far larger cohort could
resolve one. The synthetic arm behaved the same way, the effect converging
to $\dAUC\approx+0.003$ across up to $100{,}000$ networks.

How SFI was computed does not explain the null. The exact per-region $\Delta\lam_2$
SFI, which carries no first-order approximation error, also added nothing over
the competitor set. Measuring the mechanism directly, without reference to the
label, placed it partly in feature redundancy: the leading canonical correlation
between the nine SFI columns and the six competitors was $0.996$, against
$0.289$ for a permuted control, and the two aggregate summaries were
reconstructible from the competitors at $R^2=0.970$. The region maximum, the
region spread and the edge-level columns retained most of their variance
($R^2<0.25$), so they carry variance the competitors do not, and that variance
does not predict the label.

Independently of the predictive question, the validity radius
$\rhoval=\lVert\Delta L\rVert/(\lam_3-\lam_2)$ is a label-free precondition that
any Fiedler-based fragility marker can be screened against before an outcome is
seen. First-order relative error scales as $\Theta(\rhoval)$ and crosses $10\%$
at $\rhoval^\ast\approx3$, while real atria sit at median $\rhoval\approx2422$,
every mesh past the threshold (Figure~2). The spectral gap collapses as
$N^{-1.00}$ on the weighted conduction Laplacian, matching the Weyl 2-manifold
exponent to within $0.5\%$, so finer meshing degrades the single-vector
estimator in a predictable direction. The screen is portable to spectral
fragility markers used outside cardiology. It bounds estimator accuracy only,
and says nothing about how well a feature ranks cases, which must still be
tested against competitors.

The methodological result is that ordinary practice would have reported this
biomarker as working. Removing the shape-family grouping inflated AUC by
$+0.117$ ($0.756$ under random cross-validation against $0.640$ grouped), and
removing the effect-size gate allowed a $\dAUC$ of $+0.003$ to pass at
$p\approx2\times10^{-24}$ on $55{,}000$ networks (Figure~3). Both failure modes
occurred in this project's own pipeline. In silico, a $\dAUC$ of $+0.103$ at
$p=0.009$ was withdrawn once $38$ further released meshes were added. On real
patients, the clinical endpoint returned $\dAUC=+0.197$ at DeLong $p=0.0061$,
clearing both of its thresholds. The analysis plan, frozen before the outcomes
were joined, identified that increment as an artefact of baseline degradation:
the baseline was anti-predictive at AUC $0.250$ and reached only $0.447$ with SFI
added. Nine columns of Gaussian noise then reproduced it, gaining $+0.137$ on
average with $6$ of $20$ draws matching or exceeding SFI (empirical $p=0.333$;
Table~2), and the fixed-sequence stopping rule halted testing. The guard that
caught the false positive was written before the outcome column was opened, not
after the result was seen.

\subsection*{Limitations}

The in-silico labels are a monodomain simulator's reentry-inducibility verdict,
not clinical atrial fibrillation; monodomain Mitchell--Schaeffer was substituted
for the pre-registered openCARP ground truth and logged as a deviation. That
labeller is strongly resolution-sensitive and does not converge in the tested
window. The inducibility rate falls from $58\%$ at $1500$ nodes to $17\%$ at
$3000$, then rises to $37.5\%$ at $\sim\!50$k, against $34.0\%$ at the
$2000$-node operating resolution, and only $7/24$ subjects are consistent
across all six tiers. Per-subject label noise is non-differential and biases every
reported increment toward zero, so the null is partly a statement about the
labeller. By construction the clinical arm is underpowered: with $48/82$ events,
power at the $0.05$ gate is $0.24$, the minimum detectable effect is
$\dAUC\approx0.10$ to $0.15$, and $80\%$ power at $0.05$ would require $371$ to
$859$ patients, so its null is not evidence of no effect. Neither public deposit
released any demographic or clinical variable; both ship anatomical fields only.
No sociodemographic characterisation, subgroup analysis, or assessment of
whether discrimination differs across subgroups is therefore possible, and the
representativeness of these patients relative to any ablation population is
unknown. There was no patient or public involvement, this being a secondary
analysis of de-identified deposits. Each cohort derives from a single
originating centre. openCARP was run only as a robustness check and is barred
from any endpoint, its calibration gate passing on one cohort and not the other.

\subsection*{Conclusions}

A closed-form spectral fragility index did not add practically meaningful
reentry-inducibility prediction beyond existing features, at synthetic scale or
on $182$ patients of real anatomy. The single-vector estimator is additionally
ill-conditioned on real tissue and lawfully more so as meshes are refined. The
clinical endpoint was not met, and it could not have resolved an effect of the
size that would matter. We report a bounded negative and a portable, label-free
validity screen.
